% Hafta 9 — Karşı tarafın araçları
% Derleme: py -3.12 tools/latex/derle.py docs/week-9/assets/h09-09-deobfuscation.tex
\documentclass[tikz,border=8pt]{standalone}
\usepackage{cen429-cizim}
\begin{document}
\begin{tikzpicture}[node distance=10mm and 10mm]
  \node[baslik] (b) at (0,0) {Karşı tarafın araçları};
  \node[altbaslik, text width=170mm] at (b.south west) {dayanıklılık, bu araçlara karşı ölçülür};

  \node[koyukutu=32mm, below=10mm of b.south west, anchor=north west] (k1) {\kb{Gizlenmiş ikili}};
  \node[uyarikutu=36mm, right=of k1] (k2) {\kb{Disassembler /}\\\kb{decompiler}\\yapıyı geri kurmaya çalışır};
  \node[kotukutu=40mm, right=of k2] (k3) {\kb{Sembolik yürütme}\\(KLEE)\\yolları otomatik çözer\\opak yüklemi kırabilir};
  \node[iyikutu=38mm, right=of k3] (k4) {\kb{Dayanıklı yüklem}\\girdi bağımlı, yol patlatan\\$\rightarrow$ araç tıkanır};
  \draw[ok, draw=soluk] (k1.east) -- (k2.west);
  \draw[ok, draw=soluk] (k2.east) -- (k3.west);
  \draw[ok, draw=soluk] (k3.east) -- (k4.west);

  \node[dip, text width=170mm, below=10mm of k1.south -| k1, anchor=north west] {%
    Klasik kalıplar kütüphanelerde tanınır $\rightarrow$ tohumla çeşitlendirin (14. hafta).};
\end{tikzpicture}
\end{document}
